% TOP_PROBLEM: {"id":30007585,"problem_number":"LOCAL-30007585","title":"Stability of social-preference measures","category_id":21,"category":{"id":21,"name":"economics","display_name":"Economics","description":"Open economic theory statements, published research agendas, and proposed scoped specifications; question provenance and historical priority are qualified per record.","slug":"economics","order_index":21,"created_at":"2026-10-08T00:00:00Z"},"status":"research_question","external_url":null,"legacy_ids":[],"published":false,"statement":"Identify whether standardized social-preference measurements reliably predict behavior across populations after accounting for elicitation noise and variation in economic context.","economics":{"original_id":74,"field":"Behavioral and experimental economics","kind":"empirical","formalization_basis":"adapted_research_question","source_status":"research_agenda","priority_status":"earliest_found","priority_note":"A January 2022 full-paper agenda was verified. Earlier validations and replications exist; originality of the general measurement question is not established.","earliest_verified_reference":{"authors":["Michael Kosfeld","Zahra Sharafi"],"title":"The Preference Survey Module: New Evidence on Social Preferences from Tehran","year":2022,"url":"https://docs.iza.org/dp15006.pdf","doi":null,"locator":"IZA Discussion Paper 15006, §1, printed p. 2; §4 Conclusion, printed pp. 11–12 (PDF pp. 4 and 13–14).","evidence_summary":"The introduction identifies cross-sample prediction as open; the conclusion calls for more countries, larger samples, and comprehensive validation and development of global preference measures."},"current_reference":{"authors":["Michael Kosfeld","Zahra Sharafi"],"title":"The Preference Survey Module: evidence on social preferences from Tehran","year":2024,"url":"https://link.springer.com/article/10.1007/s40881-023-00151-5","doi":"10.1007/s40881-023-00151-5","locator":"§1 Introduction, paragraphs 2–4; §4 Conclusion. Published online October 19, 2023; issue June 2024, pp. 152–164.","evidence_summary":"The final article explicitly retains the question of prediction outside the original validation sample and reports mixed replication results. The original catalogue incorrectly dated this article 2025."},"formalization_note":"The cross-population validation agenda is source-supported. The risk functions, repeated-measure design, and calibration criteria operationalize it rather than reproduce a published formal conjecture."},"statusQualification":"Published question or research agenda verified at the cited locator. The mathematical specification is adapted for this catalogue and includes additional assumptions; source publication does not establish first-ever priority or certify this exact model remains unsolved. The cross-population validation agenda is source-supported. The risk functions, repeated-measure design, and calibration criteria operationalize it rather than reproduce a published formal conjecture.","statusReviewedAt":"2026-10-08"}
% FIRST_POSTED: 2026-10-08
% ENTRY_KIND: statement_only
% !TeX program = lualatex
\documentclass[11pt]{article}
\usepackage{fontspec}
\usepackage[a4paper,margin=25mm]{geometry}
\usepackage{amsmath,amssymb,mathtools}
\usepackage{xurl}
\usepackage[hidelinks]{hyperref}
\newcommand{\sref}[2]{\hyperlink{source:#1:#2}{[S#2]}}
\setlength{\emergencystretch}{3em}
\begin{document}
% problemId:30007585
\section*{Stability of social-preference measures}\label{30007585}
\subsection{Definitions and mathematical statement}
Let \(c\) index countries or participant populations and \(k\) index trust, altruism, positive reciprocity, or negative reciprocity. For individual \(i\), observe a survey vector \(S_{ic}\), predetermined demographic covariates \(X_{ic}\), and an incentivized behavioral outcome \(Y_{ick}\) in a precisely specified game. A predictor \(f_k\), trained in a source population \(c_0\), maps survey responses and covariates to behavior. For a common loss function \(\ell_k\), define its target-population error and improvement over a covariate-only predictor \(g_k\) by
\[R_{ck}=\mathbb E_c[\ell_k(Y_{ick},f_k(S_{ic},X_{ic}))],\qquad V_{ck}=\mathbb E_c[\ell_k(Y_{ick},g_k(X_{ic}))]-R_{ck}.\]
Determine whether \(V_{ck}\) is economically meaningful across populations, and whether item rankings and conditional calibration remain stable after translation, stake normalization, and changes in social context. Estimate reliability with repeated surveys and games, separating measurement noise from population changes in conditional behavior. Assess quantitative hypothetical-game items separately from qualitative self-assessments; neither a significant source-sample correlation nor a country-average association establishes individual portability. Prespecify common coding and game protocols, quantify sampling uncertainty, and test predictors on independent target samples. The research target is a validated domain of prediction and a measurement design that improves it, rather than proof that social preferences are invariant across cultures or that any single experimental game reveals an unchanging latent trait.

\par\textbf{Formulation and scope.} The cross-population validation agenda is source-supported. The risk functions, repeated-measure design, and calibration criteria operationalize it rather than reproduce a published formal conjecture.

\par\textbf{Open-question provenance.} An explicit question or research agenda was verified at the source locator. This does not by itself certify that every later version remains unresolved \sref{30007585}{1}.

\par\textbf{Historical priority.} A January 2022 full-paper agenda was verified. Earlier validations and replications exist; originality of the general measurement question is not established.

\subsection{Short English statement}
Identify whether standardized social-preference measurements reliably predict behavior across populations after accounting for elicitation noise and variation in economic context.

\subsection{Sources}
\begin{itemize}\raggedright
\item[\textnormal{[S1]}]\hypertarget{source:30007585:1}{} Michael Kosfeld, Zahra Sharafi. 2022. The Preference Survey Module: New Evidence on Social Preferences from Tehran. IZA Discussion Paper 15006, §1, printed p. 2; §4 Conclusion, printed pp. 11–12 (PDF pp. 4 and 13–14)..
\url{https://docs.iza.org/dp15006.pdf}.
{\small\textit{Earliest verified question/agenda reference; historical priority qualified below. The introduction identifies cross-sample prediction as open; the conclusion calls for more countries, larger samples, and comprehensive validation and development of global preference measures.}}

\item[\textnormal{[S2]}]\hypertarget{source:30007585:2}{} Michael Kosfeld, Zahra Sharafi. 2024. The Preference Survey Module: evidence on social preferences from Tehran. §1 Introduction, paragraphs 2–4; §4 Conclusion. Published online October 19, 2023; issue June 2024, pp. 152–164..
\url{https://link.springer.com/article/10.1007/s40881-023-00151-5}.
DOI: 10.1007/s40881-023-00151-5.
{\small\textit{Supporting or current literature; see evidence qualification. The final article explicitly retains the question of prediction outside the original validation sample and reports mixed replication results. The original catalogue incorrectly dated this article 2025.}}
\end{itemize}
\end{document}
